\section{Pre-registration record and deviations}\label{app:prereg}

\Cref{tab:runs} lists the runs reported in the paper. Four of them decide its claims, and
\cref{tab:claims,tab:holdout} give their pre-registered claims with the outcomes;
\cref{tab:criteria} lists every criterion fixed in advance for every run, with the value
measured and the outcome as recorded by the code at the time. The pre-registration of the
July run also retired, before its run, an earlier criterion (a displacement improvement of
at least 10\% from the energy term with 64 labels) that an earlier run had already answered
in the negative; the July run measured \numGOneRetired{} with the fixed-weight form of the
energy term, and the criterion played no part in its decisions. The following departures
from the plans bear on the results. They are recorded in the repository's list of deviations
and its provenance note, except the remark at the end of the paragraph on the energy term,
which was noticed while the paper was written, and the paragraphs on the two-dimensional
comparison, the holdout evaluation and the three-dimensional energy-norm run, whose content
is recorded with those runs' pre-registrations and records.

\begingroup\footnotesize\setlength{\tabcolsep}{3pt}\setlength{\LTcapwidth}{\linewidth}
\begin{longtable}{@{}>{\raggedright\arraybackslash}p{0.21\linewidth}
>{\raggedright\arraybackslash}p{0.14\linewidth}
>{\raggedright\arraybackslash}p{0.17\linewidth}
>{\raggedright\arraybackslash}p{0.275\linewidth}
>{\raggedright\arraybackslash}p{0.135\linewidth}@{}}
\caption{The runs reported in the paper, by the names used in the text. The runs of July
2026 used the earlier code, the others the later code, which multiplies the decoded field by
$F_{\max}$ (\cref{sec:methods:networks}). Section~\ref{app:repro} lists the configurations
and records with their hashes.}\label{tab:runs}\\
\toprule
Run & Date & Pre-registered & Role & Reported in \\
\midrule
\endfirsthead
\caption[]{The runs reported in the paper (continued)}\\
\toprule
Run & Date & Pre-registered & Role & Reported in \\
\midrule
\endhead
\bottomrule
\endfoot
Three-dimensional comparison & 22--28 September 2026 & yes, as an amendment & the label-free,
displacement-supervised and graph networks in three dimensions; transfer to a finer mesh &
\cref{sec:labelfree:three,sec:checks}, \ref{supp:three}, \ref{sec:transfer} \\
Initial three-dimensional run & 31 August to 21 September 2026 & yes & superseded: its label-free
objective was defective (deviation D14) & \ref{app:prereg} \\
Two-dimensional comparison & 8--9 October 2026 & yes & the norm of the loss against the
labels (H1--H3) & \cref{sec:norm:two,sec:labelfree:two}, \ref{supp:two} \\
Holdout evaluation & 21 October 2026 & yes & the claims of both comparisons on instances
that no network had seen; no training & \cref{sec:holdout} \\
Three-dimensional energy-norm run & 22--25 October 2026 & yes & the norm of the loss
against the labels in three dimensions & \cref{sec:norm:three}, \ref{supp:three} \\
Ritz baseline & 28 October 2026 & no; specified before any paired computation & a classical
approximation, paired with every network & \cref{sec:cost:ritz}, \ref{sec:transfer},
\ref{supp:ritz} \\
Latent-regularisation run (2D) & 29--30 September 2026 & yes & the two-dimensional
label-free networks; latent regularisation & \cref{sec:labelfree:two}, \ref{app:e1} \\
Bottleneck run (3D) & 30 September 2026 & yes & a token bottleneck to reduce the cost of
attention & \cref{sec:cost:labels}, \ref{app:e1} \\
July run (2D) & 16 July 2026 & yes; document uploaded with its hash line blank & the earlier
two-dimensional comparison; warm start (K5); cross-resolution term (K6) &
\ref{supp:warm}, \ref{app:july} \\
Diagnostic run (2D) & 30 July 2026 & no & replication of the label-free network; the energy
term with 16 labels & \ref{app:prereg}, \ref{app:july} \\
Joint-embedding run (2D) & 31 July 2026 & yes; document committed after the run & the
joint-embedding term (K3); the scope of the paper & \ref{app:prereg} \\
\end{longtable}
\endgroup

% generated from the committed records; do not edit
\begingroup\small\setlength{\LTcapwidth}{\linewidth}
\begin{longtable}{@{}>{\raggedright\arraybackslash}p{0.145\linewidth}>{\raggedright\arraybackslash}p{0.405\linewidth}>{\raggedright\arraybackslash}p{0.21\linewidth}>{\raggedright\arraybackslash}p{0.12\linewidth}@{}}
\caption{Every pre-registered criterion of the runs reported here, with its outcome.}\label{tab:criteria}\\
\toprule
Run & Criterion & Measured & Outcome \\
\midrule
\endfirsthead
\caption[]{Every pre-registered criterion of the runs reported here, with its outcome. (continued)}\\
\toprule
Run & Criterion & Measured & Outcome \\
\midrule
\endhead
\bottomrule
\endfoot
July run (2D), 16 July 2026 & G1$'$ (a), energy-term sub-experiment (separately trained networks): supervised + energy term at least 3 times better than the zero field in displacement, and better than the two naive baselines and the plain polynomial, fitted on 1,024 labelled instances, at every budget & 4.5--6.1 times; baselines beaten at every budget & met \\
 & G1$'$ (b), \cref{tab:2djuly,tab:labeleffjuly}, 64 labels: the energy term lowers the relative energy gap by $\geq$ 50\% and the von Mises error by $\geq$ 25\% & 78.7\%, 65.4\% & met \\
 & G1$'$ (c), \cref{tab:2djuly,tab:labeleffjuly}, 64 labels: label-free, then supervised beats supervised by $\geq$ 10\% in relative energy gap or $\geq$ 5\% in displacement & 92.6\%, 2.0\% & met \\
 & Gate G1$'$ = (a) and (b) and (c) &  & GO \\
 & K1, energy-term sub-experiment: the better of the fixed and the gradient-scaled energy term lowers the relative energy gap by $<$ 25\% at every budget & 98.0\% / 94.4\% / 90.6\% / 74.0\% & not triggered \\
 & K2, \cref{tab:2djuly}: label-free displacement error $>$ 30\% above supervised, 1,024 labels & +3.6\% & not triggered \\
 & C1-advantage, \cref{tab:labeleffjuly}: label-free relative energy-gap advantage over supervised $<$ 40\% at every budget & 98.7\% / 97.0\% / 93.5\% / 80.4\% & not triggered \\
 & K4: a latent-space regulariser raises the standardised effective rank of the latent vectors by $\leq$ 1.5 times & 1.70 times & not triggered \\
 & K5: conjugate gradients from the prediction of a label-free network (64 validation instances) save $<$ 20\% of the iterations from zero & $-$0.9\% & triggered \\
 & K6: a cross-resolution invariance term shrinks the transfer gap by $<$ 10\% at coarsening 2.5 & $-$23.7\% & triggered \\
 & E6 (alignment): within-geometry rank correlation of latent and solution distances $<$ 0.3 & 0.730 & not triggered \\
 & C5: any pre-registered check of the conditioning bound, the modewise identity or the conjugate-gradient bound violated & none & not triggered \\
\midrule
Joint-embedding run (2D), 31 July 2026 & K3: a self-supervised joint-embedding term added to label-free pretraining changes the displacement error and the relative energy gap of the transformer, fine-tuned with 16, 64 and 256 labels, by less than 3\% at every budget (if triggered, the term is dropped) & change with the term, 16 / 64 / 256 labels: displacement $-$13.5\% / $-$0.5\% / +2.6\%; relative energy gap $-$48.4\% / +92.5\% / $-$16.4\% (positive: more accurate) & not triggered \\
\midrule
Initial three-dimensional run, 21 September 2026 & Gate G2 and kill conditions as registered; the label-free objective was defective (deviation D14) & (a), (b), (c) not met; KP1, KP2, KP4 triggered & NO-GO \\
Pilot of the amendment (3D), 22 September 2026 & Label-free displacement error after 20 epochs $<$ 0.90 & 0.309 & met \\
Three-dimensional comparison, 28 September 2026 & G2 (a), budgets $\geq$ 64 (amendment), \cref{tab:labeleff}: displacement-supervised transformer with the energy term at least 3 times better than the zero field in displacement, and better than both naive baselines & 3.9 / 12.0 / 21.3 times; naive beaten & met \\
 & G2 (a), every budget (the registered form, reported for reference) & 2.71 times at 16 labels & not met \\
 & G2 (b), \cref{tab:labeleff,tab:3d}: label-free displacement error within +10\% of the displacement-supervised transformer's at 1,024 labels, and relative energy-gap advantage $\geq$ 40\% at every budget & $-$36.3\%; 99.95\% / 99.85\% / 98.86\% / 99.97\% & met \\
 & G2 (c), \cref{tab:transfer}: finer mesh, zero-shot displacement error $\leq$ 1.25 times in-band, and better than both naive baselines & 8.62 times; best naive 0.401 against 0.258 & not met \\
 & Gate G2 = (a) and ((b) or (c)), amended &  & GO \\
 & Gate G2, registered form (reference) &  & NO-GO \\
 & KP1: supervised (1,024 labels) better than label-free by $>$ 10\% in displacement & $-$36.3\% & not triggered \\
 & KP2: relative energy-gap advantage $<$ 40\% at any budget & minimum 98.9\% & not triggered \\
 & KP3: the energy term lowers the relative energy gap by $<$ 25\% at every budget & +96.6\% / +85.7\% / $-$20.7\% / +99.6\% & not triggered \\
 & KP4: finer mesh, zero-shot displacement error $>$ 1.5 times in-band, or a naive baseline better & 8.62 times & triggered \\
 & KP5: any pre-registered check of the conditioning bound, the modewise identity or the conjugate-gradient bound violated & none & not triggered \\
 & KP6 (alignment): within-geometry rank correlation of latent and solution distances $<$ 0.3 & 0.827 & not triggered \\
\midrule
Latent-\hspace{0pt}regularisation run (2D), 29 September 2026 & Latent regularisation, K1: relative change of the displacement error or of the relative energy gap beyond the noise guard & within the guard & not triggered \\
 & Latent regularisation, K2: the separation does not rise in any seed & rose in every seed & not triggered \\
 & Latent regularisation, GO: separation raised by $\geq$ 0.02 in every seed & +0.009 / +0.014 / +0.002 & NO-GO \\
\midrule
Bottleneck run (3D), 30 September 2026 & Token bottleneck, M = 512, K1: in-band relative energy gap or finer-mesh displacement error worse than the label-free transformer's beyond the noise guard & +280.9\% (guard 287.2\%), +79.1\% (guard 41.9\%) & triggered \\
 & Token bottleneck, M = 512, K2: finer-mesh training step longer than 2 s & 0.049 s & not triggered \\
 & Token bottleneck, M = 512: outcome &  & dropped \\
 & Token bottleneck, M = 1,024, K1: in-band relative energy gap or finer-mesh displacement error worse than the label-free transformer's beyond the noise guard & +183.3\% (guard 70.5\%), +122.0\% (guard 33.3\%) & triggered \\
 & Token bottleneck, M = 1,024, K2: finer-mesh training step longer than 2 s & 0.047 s & not triggered \\
 & Token bottleneck, M = 1,024: outcome &  & dropped \\
\midrule
Two-dimensional comparison, 9 October 2026 & H1, \cref{tab:2d,tab:cm2d}: label-free transformer lower than the displacement-supervised transformer ($\mathcal{L}_D$) in relative energy gap beyond the noise guard, 1,024 labels & $-$71.3\% (guard 12.4\%) & supported \\
 & H2a: energy-norm-supervised transformer ($\mathcal{L}_K$) lower than the displacement-supervised transformer ($\mathcal{L}_D$) in relative energy gap beyond the guard & $-$78.2\% (guard 16.4\%) & supported \\
 & H2b: energy-norm-supervised transformer lower than the displacement-supervised transformer in von Mises error beyond the guard & $-$54.9\% (guard 12.3\%) & supported \\
 & H3 (a reading, no criterion): energy-norm-supervised transformer against the label-free transformer in relative energy gap & $-$24.2\% (guard 38.1\%) & no difference shown \\
 & Reuse: the label-free networks of the latent-regularisation run reproduce that run's per-instance relative energy gaps and displacement errors (largest relative deviation at most $10^{-4}$) & largest relative deviation 0 & passed \\
\midrule
Holdout evaluation, 21 October 2026 & T1 (3D): label-free transformer better than the displacement-supervised transformer on all five metrics, in each seed and in the seed mean & 20 of 20 comparisons & replicated \\
 & T2 (3D): graph network's displacement error lower than the label-free transformer's on more than 50\% of the instances in each seed, and its von Mises error higher on at least 95\% of the instance-seed pairs & 86.3\%--95.8\%; 100\% & replicated \\
 & T3 (3D): each displacement-supervised network has load cases worse than the zero field, each with a positive energy and no other & 340 and 237 of 8,232 load cases & replicated \\
 & H1 (2D): label-free transformer lower than the displacement-supervised transformer ($\mathcal{L}_D$) in relative energy gap beyond the noise guard & $-$71.3\% (guard 11.5\%) & replicated \\
 & H2a (2D): energy-norm-supervised transformer ($\mathcal{L}_K$) lower than the displacement-supervised transformer in relative energy gap beyond the guard & $-$42.9\% (guard 47.4\%) & not replicated \\
 & H2b (2D): energy-norm-supervised transformer lower than the displacement-supervised transformer in von Mises error beyond the guard & $-$52.8\% (guard 12.9\%) & replicated \\
 & H2c (2D): energy-norm-supervised transformer's von Mises error less than half the displacement-supervised transformer's, rel below $-$50\% & $-$52.8\% & replicated \\
 & D2 (2D): label-free transformer's displacement error higher than the displacement-supervised transformer's and its other four metrics better, in each seed and in the seed mean & 20 of 20 comparisons & replicated \\
 & H3 (2D, a reading without a criterion): energy-norm-supervised transformer against the label-free transformer in relative energy gap & +99.1\% (guard 160.2\%) & no difference shown \\
 & V: the per-instance values of every network and metric distributed alike on the holdout and the validation instances (two-sided Mann--Whitney tests, Holm's procedure at a family-wise level of 5\%) & 35 quantities, zero rejected & \hspace{0pt}representative \\
\midrule
Three-dimensional energy-norm run, 25 October 2026 & H2a-3D: energy-norm-supervised transformer ($\mathcal{L}_K$) lower than the displacement-supervised transformer ($\mathcal{L}_D$) in relative energy gap beyond the noise guard, on the seeds' medians & $-$98.2\% (guard 11.7\%) & supported \\
 & H2b-3D: the same in von Mises error & $-$83.6\% (guard 10.0\%) & supported \\
 & H3-3D (a reading without a criterion): energy-norm-supervised transformer against the label-free transformer in relative energy gap, on the seeds' medians & $-$41.0\% (guard 27.0\%) & energy-norm-supervised transformer lower \\
\end{longtable}
\endgroup
\vspace{-\smallskipamount}
\noindent\begin{minipage}{\linewidth}\footnotesize\raggedright
Runs are named as in \cref{tab:runs} and dated by their records; the pilot of the amendment is the short run that the amendment required before the three-dimensional comparison (Section~\ref{app:prereg}). Labels are those of the pre-registration documents, each of which numbers its own criteria. Gates (G1$'$, G2) decide whether the study proceeds to its next stage. Kill conditions withdraw a claim, or abandon a component, when triggered: in two dimensions K1--K6, C1-advantage (the claimed advantage in relative energy gap), C5 (the pre-registered checks of Section~\ref{app:repro:checks}) and E6 (latent alignment); in three dimensions KP1--KP6, with KP5 the counterpart of C5. The latent-regularisation and bottleneck runs number their criteria afresh, with the meanings given here. A criterion reads the trained networks of the tables it names; the energy-term sub-experiment of the July run trained its own supervised networks with and without the energy term, with the same configuration. K4, E6 and KP6 concern latent representations, which this paper does not use. Standardised effective rank: the participation ratio of the eigenvalues of the covariance of the latent vectors, each dimension standardised. The 3D amendment, the deviations and the joint-embedding term are described below, the latent regularisation and the token bottleneck in Section~\ref{app:e1}, the July run in Section~\ref{app:july}.\par
Noise guard: a comparison of a network A with a reference B over three seeds each, with rel = mean(A)/mean(B) $-$ 1 and SE$_\mathrm{rel}$ = $(s_A^2/3 + s_B^2/3)^{1/2}$ divided by mean(B), where $s_A$ and $s_B$ are the sample standard deviations over the seeds, finds A lower beyond the guard when rel is below minus the guard, max(10\%, 2 SE$_\mathrm{rel}$), and higher beyond it when rel exceeds the guard; within it no difference is shown, which is not equivalence. The hypotheses H1--H3 of the 2D comparison use it (\cref{tab:cm2d}), and so do the latent-regularisation and bottleneck runs, with the label-free network as B (their pre-registrations call it the resolution), the holdout evaluation, and the three-dimensional energy-norm run, which applies it to each seed's median over the instances (\cref{tab:k3d}). $\mathcal{L}_D$ and $\mathcal{L}_K$ are the losses \eqref{eq:ld} and \eqref{eq:lk}.
\end{minipage}
\par\medskip


\paragraph{Deviation D14}
The initial three-dimensional run (31 August to 21 September 2026) needed eight attempts
because of memory and process faults; each restart reused the networks already trained under
the same registered configuration, and the engineering changes the restarts needed are
recorded as leaving every value unchanged. It completed its registered protocol and read
NO-GO, with conditions (a), (b) and (c) of its gate unmet and kill conditions KP1, KP2 and
KP4 triggered. After the verdict, its label-free networks were found to have a displacement
error of \numDfourteenDisp{} and a relative energy gap of \numDfourteenGap{}, identical to
five decimals across seeds and consistent with predicting the solution scaled by
$\alpha=\numDfourteenAlpha{}$ (displacement error $1-\alpha$, relative energy gap
$(1-\alpha)^2$). The cause was in the code: inference multiplied the decoded field by
$F_{\max}$, the label-free objective scored it without that factor, and on this corpus
$F_{\max}\approx\alpha$. The fix gives the network one decoding path, used by the objective
and by inference, with a test that fails on the registered code. The supervised networks,
which are trained and evaluated through the inference path, were not affected. The verdict
of the initial run stands as recorded. A pre-registered amendment then repeated everything
that depended on the label-free objective: the label-free networks, their fine-tuning on the
finer mesh and the latent-alignment criterion read from them. It reused the 30 trained
supervised networks and the 6 networks trained from scratch on the fine mesh, each checked
against its recorded SHA-256 hash, and it required a short pilot run first: one seed, 20
epochs, which had to bring the displacement error below \numPilotLimit{} and reached
\numPilotDisp{}. A first version of the amendment was registered and withdrawn on the same
day, before any run, because the audit of its execution path had not been completed; the
second version was registered after the audit. The amended run, the three-dimensional
comparison (22 to 28 September 2026), was executed once, without interruption.

\paragraph{The sanity condition of the three-dimensional gate}
Condition (a) of the three-dimensional gate required the supervised transformer with the
energy term to be at least three times more accurate in displacement than the zero field,
and better than both naive baselines, at every budget. In the initial run it failed at 16
labels (\numGTwoSanitySixteen{} times), independently of D14. The amendment, written with
this result known, assessed condition (a) only from the decision budget of 64 labels, where
it was met by factors of \numGTwoSanitySixtyFour{} to \numGTwoSanityMax{}. Its stated
rationale is that the factor of three had been calibrated on the two-dimensional family,
and that on the three-dimensional family the supervised network itself sits below it with
16 labels and above it from 64 labels on. In two dimensions the condition had read the
networks of the energy-term sub-experiment of the July run, which met it at every budget;
the networks of \cref{tab:labeleffjuly} would not have met it with 16 labels
(\numGOneSanityTableSixteen{} times), because of the scatter described in
Section~\ref{app:july}.
The amendment declared the change as made after the fact and required the original form to
be computed and reported beside it. Under the amendment the gate read GO; under the
registered form, which assesses every budget, it read NO-GO. Both readings are in
\cref{tab:criteria}. The change affects only the gate's verdict, not any measurement
reported in this paper.

\paragraph{Deviation D16}
The text of the three-dimensional pre-registration states that the label-free network
trains on the 1,744 in-band instances outside the validation set. The registered
configuration, which is what ran, trained it on the first 1,024 of them, the same instances
on which the supervised networks were trained with labels. Every three-dimensional run
executed this way. The discrepancy was found after the amended run and is recorded with the
other deviations. It works against the label-free network, which saw fewer instances than
the text promised, and it leaves the cleaner comparison: the same instances, with and
without their solutions.

\paragraph{The energy term with 16 labels}
In the networks of \cref{tab:2djuly,tab:labeleffjuly}, the two-dimensional supervised
transformer with the energy term used the gradient-scaled form at every budget, and with 16
labels its displacement errors scattered across seeds (Section~\ref{app:july}); the
energy-term sub-experiment of the same run trained both the gradient-scaled and the
fixed-weight form. The diagnostic run (30 July 2026), with six seeds in two separate
experiments, gave displacement errors from \numDiagScaledMin{} to \numDiagScaledMax{} for
the gradient-scaled form, against \numDiagFixedMin{} to \numDiagFixedMax{} in all
\numDiagFixedSeeds{} seeds for a fixed weight of one on the normalised energy. The
three-dimensional pre-registration therefore fixed that weight at 16 labels and the
gradient-scaled form from 64 on. The description of the protocol stored in the
three-dimensional report still names the gradient-scaled form for every budget; the
registered code applied the pre-registered policy.

\paragraph{The joint-embedding term (criterion K3)}
Criterion K3 of the July run's pre-registration concerned a separate run, the
joint-embedding run (31 July 2026). That run pretrained the transformer without labels for
200 epochs, once on the energy alone and once on the energy together with a self-supervised
joint-embedding term, which predicts the latent vectors of masked regions of the mesh from
those of the rest, with a regulariser on the mean latent vectors; it then fine-tuned both
with 16, 64 and 256 labels for 200 epochs, with three seeds. K3 would have dropped the term,
and with it the scope that the pre-registration planned for the paper, had the term changed
neither the displacement error nor the relative energy gap by more than \numKthreeBand{} at
any budget. K3 was not triggered, and the registered rule therefore kept the term
(\cref{tab:criteria}). The criterion did not state a direction. Of the \numKthreeOutside{}
changes outside the band, \numKthreeWorse{} made the fine-tuned network less accurate; the
remaining one, a relative energy gap \numKthreeGain{} lower with 64 labels, compares with a
baseline whose seeds scattered widely (relative energy gaps of \numKthreeBaseSeeds{},
against \numKthreeBaseRef{} for the same configuration in the July run). The paper reports
the energy objective without the term. That scope was chosen after the result was known, and
it departs from the scope that the rule selected.

\paragraph{The two-dimensional comparison}
The two-dimensional comparison was registered on 8 October 2026 and run once, starting on 8
October and completing on 9 October 2026; its adjudication, by a script fixed with the
registration, recorded no deviation, and two operational departures recorded with the run,
the route by which its return was fetched and a log file renamed after the run, touched no
file that the adjudication reads. \Cref{tab:cm2d} gives its hypotheses with the readings
that the pre-registration required beside each verdict. Its label-free row reused the three
label-free networks of the latent-regularisation run (Section~\ref{app:e1}): each was
checked against the SHA-256 hash that run had recorded, and evaluated again they reproduced
that run's per-instance validation values exactly. Three things were known when the
comparison was registered: the label-free networks' values on the validation instances,
recorded with the latent-regularisation run; the results of the July run that motivated the
comparison, measured on the same validation instances, so that the comparison has no fresh
evaluation set; and the outcome of a functional check of the energy-norm loss, run before
the registration was stamped, which had trained a transformer of width 32 and depth 2 for 25
epochs on 48 instances of another two-dimensional corpus, with two seeds. At that far
smaller scale the check had shown both reductions that H2a and H2b test, and a label-free
relative energy gap below that of supervised training on the displacement error, the
comparison of H1, with the energy-norm loss and the label-free objective ranked differently
in its two seeds, the comparison of H3. The registration disclosed the check. The readings
of \cref{sec:norm:spectrum}, and the rule that chose the instance of \cref{fig:field2d},
were specified after the verdict and bear on no criterion.

% generated from the committed records; do not edit
\begin{table}[tbp]
\centering
\footnotesize
\setlength{\tabcolsep}{4pt}
\caption{The pre-registered hypotheses of the two-dimensional comparison, with 1,024 labelled instances, as its adjudicator recorded them, and the readings its pre-registration required beside each.}
\label{tab:cm2d}
\resizebox{\ifdim\width>\textwidth\textwidth\else\width\fi}{!}{%
\begin{tabular}{lcccc}
\toprule
 & H1 & H2a & H2b & H3 \\
\midrule
Network A & Label-free & Supervised, energy norm & Supervised, energy norm & Supervised, energy norm \\
Reference B & Supervised, displacement & Supervised, displacement & Supervised, displacement & Label-free \\
Metric & Relative energy gap & Relative energy gap & von Mises error & Relative energy gap \\
\midrule
A, seeds 0, 1, 2 & 0.0374, 0.0405, 0.0371 & 0.0370, 0.0146, 0.0355 & 0.119, 0.0829, 0.116 & 0.0370, 0.0146, 0.0355 \\
B, seeds 0, 1, 2 & 0.132, 0.148, 0.120 & 0.132, 0.148, 0.120 & 0.232, 0.251, 0.222 & 0.0374, 0.0405, 0.0371 \\
\midrule
rel & $-$71.3\% & $-$78.2\% & $-$54.9\% & $-$24.2\% \\
SE$_\mathrm{rel}$ & 6.2\% & 8.2\% & 6.1\% & 19.0\% \\
Guard & 12.4\% & 16.4\% & 12.3\% & 38.1\% \\
Reading & supported & supported & supported & no difference shown \\
\midrule
Medians: rel (guard) & $-$78.7\% (10.0\%) & $-$83.3\% (11.6\%) & $-$53.6\% (11.8\%) & $-$21.4\% (33.7\%) \\
Welch 95\% interval & $-$97\% to $-$45\% & $-$101\% to $-$55\% & $-$73\% to $-$37\% & $-$103\% to +54\% \\
Instance resampling 95\% & $-$74.3\% to $-$68.3\% & $-$80.3\% to $-$76.3\% & $-$57.2\% to $-$52.8\% & $-$28.0\% to $-$21.9\% \\
\bottomrule
\end{tabular}}

\smallskip
\begin{minipage}{\linewidth}\footnotesize\raggedright
rel, SE$_\mathrm{rel}$ and the noise guard as defined in the note to \cref{tab:criteria}, over the three seeds' values, each the mean over the 256 validation instances and their four load cases.\par
Medians: the same reading on each seed's median over the instances instead of its mean. Welch: rel $\pm$ t $\times$ SE$_\mathrm{rel}$, with t the 97.5\% quantile of Student's distribution at the Welch--Satterthwaite degrees of freedom, a symmetric interval that can extend below $-$100\%. Instance resampling: the 2.5\% and 97.5\% percentiles of rel over 2,000 resamplings of the 256 validation instances, the same instances drawn for both networks and every seed, the seeds kept.\par
H3 is a reading without a criterion. With the roles of A and B exchanged, rel was +31.9\% against a guard of 50.2\%.
\end{minipage}
\end{table}


\paragraph{The holdout evaluation}
The holdout evaluation was registered and stamped on 18 October 2026 and run on 21 October
2026, before the three-dimensional energy-norm run; its adjudicator, fixed with the
registration, read the eight claims, H3 and V (\cref{tab:holdout}). Its sets were fixed with
it. In three dimensions the corpus has \numHoldThreeTail{} instances beyond the validation
set and the training pool (\cref{sec:methods:data}); the holdout set is \numHoldThreeN{} of
them, without the \numHoldThreeExcl{} that other recorded sets had drawn again. In two
dimensions the corpus has \numHoldTwoTail{} instances beyond the deepest prefix of its
training pool that any recorded run read; the holdout set is \numHoldTwoN{} of them, without
the \numHoldTwoExcl{} that another recorded tool had rebuilt or whose own files were held
outside the training machine. The evaluation labelled them with the sparse direct solver, in
\numHoldThreeSolves{} solves in three dimensions and \numHoldTwoSolves{} in two, and
evaluated the stored networks on them; the displacement-supervised transformer with the
energy term, whose networks were not stored, was not evaluated. It trained nothing: each of
its \numHoldStates{} stored networks was checked against its recorded hash and, before the
held-out instances, reproduced its recorded per-instance values on the validation instances,
the transformers exactly and the graph networks to within \numHoldContentDev{} in the median
of the displacement error and of the relative energy gap, the two metrics that it checked.
The pre-registration required the paper to report every claim's verdict and V's, to qualify
a claim that did not replicate, with the held-out values beside the original ones, and to
report the held-out means and medians beside the validation ones;
\cref{sec:holdout,sec:norm:two} and \cref{tab:holdout,tab:holdoutvalues} do so. Two
departures were recorded with the run. Each of its four steps started twice: the first
attempts ran on a checkout of an earlier commit that did not contain the steps' script, and
each stopped at once without writing anything; the steps were then run on the stamped
commit. And between the search of the sandbox for other sets drawn from the corpora, on
which the definition of the holdout sets rests, and the run, an independent review of the
code of the three-dimensional energy-norm run, on 20 October 2026, meshed three validation
instances of the three-dimensional corpus to compare their node counts with the records, and
meshed the 1,024 instances of its training pool at their own mesh sizes, beyond the bound of
80 such instances that the pre-registration set; it rebuilt and solved two of them and ran
one training step of the full-size transformer on each; it replayed the generator's draws
for all 2,000 instances of the corpus, the held-out ones included, without meshing them, and
meshed, solved or read no held-out instance. Neither departure changed a label, an
evaluation or a criterion.

% generated from the committed records; do not edit
\begingroup\footnotesize\setlength{\LTcapwidth}{\linewidth}\setlength{\tabcolsep}{3pt}
\begin{longtable}{@{}llccc@{}}
\caption{Holdout evaluation: the mean and the median of every network and metric on the validation and on the held-out instances.}\label{tab:holdoutvalues}\\
\toprule
Network & Metric & \shortstack{Validation: \\ mean / median} & \shortstack{Holdout: \\ mean / median} & \shortstack{Ratio of the \\ means (95\%)} \\
\midrule
\endfirsthead
\caption[]{Holdout evaluation: the mean and the median of every network and metric on the validation and on the held-out instances. (continued)}\\
\toprule
Network & Metric & \shortstack{Validation: \\ mean / median} & \shortstack{Holdout: \\ mean / median} & \shortstack{Ratio of the \\ means (95\%)} \\
\midrule
\endhead
\bottomrule
\endfoot
\multicolumn{5}{@{}l}{\emph{Three dimensions, 686 held-out instances}} \\
Label-free & Displacement & 0.0300 / 0.0280 & 0.0285 / 0.0269 & 0.95 (0.93--0.98) \\
 & Energy gap & 0.00912 / 0.00758 & 0.00805 / 0.00739 & 0.88 (0.86--0.91) \\
 & von Mises & 0.0554 / 0.0532 & 0.0541 / 0.0521 & 0.98 (0.96--0.99) \\
 & Peak stress & 0.0807 / 0.0788 & 0.0799 / 0.0775 & 0.99 (0.97--1.01) \\
 & Recall & 0.793 / 0.790 & 0.790 / 0.787 & 1.00 (0.99--1.00) \\
Supervised, displacement & Displacement & 0.0470 / 0.0408 & 0.0493 / 0.0419 & 1.05 (0.98--1.16) \\
 & Energy gap & 31.2 / 0.255 & 172 / 0.231 & 5.50 (0.01--16.49) \\
 & von Mises & 0.370 / 0.259 & 0.408 / 0.253 & 1.10 (0.75--1.80) \\
 & Peak stress & 0.784 / 0.235 & 0.646 / 0.209 & 0.82 (0.37--1.71) \\
 & Recall & 0.619 / 0.629 & 0.623 / 0.635 & 1.01 (1.00--1.01) \\
Graph network & Displacement & 0.0168 / 0.0137 & 0.0160 / 0.0142 & 0.95 (0.91--1.01) \\
 & Energy gap & 0.447 / 0.209 & 0.307 / 0.197 & 0.69 (0.59--0.84) \\
 & von Mises & 0.296 / 0.252 & 0.276 / 0.242 & 0.93 (0.90--0.97) \\
 & Peak stress & 0.488 / 0.186 & 0.353 / 0.172 & 0.72 (0.63--0.82) \\
 & Recall & 0.679 / 0.700 & 0.683 / 0.702 & 1.01 (1.00--1.01) \\
\midrule
\multicolumn{5}{@{}l}{\emph{Two dimensions, 4,111 held-out instances}} \\
Label-free & Displacement & 0.0634 / 0.0472 & 0.0646 / 0.0477 & 1.02 (0.99--1.05) \\
 & Energy gap & 0.0383 / 0.0253 & 0.0393 / 0.0284 & 1.02 (0.99--1.06) \\
 & von Mises & 0.123 / 0.120 & 0.124 / 0.124 & 1.01 (0.99--1.03) \\
 & Peak stress & 0.128 / 0.138 & 0.129 / 0.140 & 1.01 (0.99--1.03) \\
 & Recall & 0.753 / 0.753 & 0.751 / 0.752 & 1.00 (0.99--1.00) \\
Supervised, displacement & Displacement & 0.0505 / 0.0406 & 0.0519 / 0.0410 & 1.03 (1.00--1.05) \\
 & Energy gap & 0.133 / 0.118 & 0.137 / 0.119 & 1.02 (0.99--1.07) \\
 & von Mises & 0.235 / 0.231 & 0.236 / 0.231 & 1.00 (0.99--1.01) \\
 & Peak stress & 0.141 / 0.147 & 0.145 / 0.147 & 1.03 (1.01--1.05) \\
 & Recall & 0.668 / 0.675 & 0.669 / 0.677 & 1.00 (1.00--1.00) \\
Supervised, energy norm & Displacement & 0.0547 / 0.0420 & 0.0567 / 0.0419 & 1.04 (1.01--1.07) \\
 & Energy gap & 0.0290 / 0.0210 & 0.0781 / 0.0223 & 2.69 (1.12--4.64) \\
 & von Mises & 0.106 / 0.108 & 0.111 / 0.109 & 1.05 (1.02--1.08) \\
 & Peak stress & 0.0947 / 0.100 & 0.105 / 0.101 & 1.11 (1.02--1.24) \\
 & Recall & 0.804 / 0.802 & 0.802 / 0.803 & 1.00 (1.00--1.00) \\
Graph network & Displacement & 0.0499 / 0.0392 & 0.0496 / 0.0383 & 1.00 (0.97--1.02) \\
 & Energy gap & 0.825 / 0.517 & 0.793 / 0.554 & 0.96 (0.92--1.00) \\
 & von Mises & 0.578 / 0.489 & 0.567 / 0.504 & 0.98 (0.96--1.00) \\
 & Peak stress & 0.321 / 0.153 & 0.320 / 0.159 & 1.00 (0.95--1.05) \\
 & Recall & 0.517 / 0.549 & 0.520 / 0.544 & 1.01 (1.00--1.01) \\
\end{longtable}
\endgroup
\vspace{-\smallskipamount}
\noindent\begin{minipage}{\linewidth}\footnotesize\raggedright
Per network and metric, the mean and the median over the instances of each instance's value averaged over the seeds; the validation values recomputed by the holdout evaluation, whose check of the reports' displacement errors and relative energy gaps they passed (the graph networks' to within $2 \times 10^{-4}$ in the median). Ratio: holdout mean over validation mean, with a 95\% percentile interval from 2,000 resamplings of the holdout instances. V's two-sided Mann--Whitney tests compare the per-instance values; Holm's procedure rejected none of their p-values, which the evaluation's verdict records, at a family-wise level of 5\%. The energy-norm-supervised transformer was evaluated in 2D only: its 3D counterpart was trained after this evaluation (\cref{sec:norm:three}).
\end{minipage}
\par\medskip


\paragraph{The three-dimensional energy-norm run}
The three-dimensional energy-norm run was registered and stamped on 21 October 2026, after
the holdout evaluation had returned and before any network had been trained with the
energy-norm loss in three dimensions (the review above had run one training step with either
loss on each of its two instances, on a CPU, to check the code), and run once, from 22 to 25
October 2026, without a deviation; its adjudicator recorded the verdicts and readings of
\cref{tab:k3d}. Its pre-registration fixed the statistic, each seed's median over the
validation instances, after the per-instance values of the displacement-supervised and the
label-free transformers, and the readings of the two-dimensional comparison on the medians,
were known, and gave the reason: one instance carries the displacement-supervised
transformer's seed means, so that on seed means the guard of its relative energy gap is
\numKdMeansGapTau{} and no reduction could be shown. The seed means are reported beside
every verdict, with their Welch intervals, as the pre-registration required, and so are the
readings on the held-out instances, which it fixed without a verdict. The run evaluated the
six stored networks again before training, and they reproduced their recorded values
exactly; it trained the energy-norm-supervised transformer with the displacement-supervised
transformer's settings, unchanged and not tuned for the new loss.

\paragraph{Provenance of the two-dimensional runs of July 2026}
The two-dimensional runs of July 2026 were executed outside a version-controlled checkout,
so
their reports record no commit. The code checked their configuration hashes against the
pre-registration documents at start-up, and the code version was established afterwards
from the runs' console output and by comparing the source trees of the released packages.
The repository was created, public, on 14 July 2026 at 03:23 UTC; the security log of the
account that created it, exported on 10 October 2026, records no change of its visibility
up to that day. The pre-registration document and the configuration of the July run were
first committed to it at 03:25 UTC
that day, by an upload through the hosting service's web interface, which created and
signed that commit (\texttt{b365af5}), the repository's first. The commit holds the code
release that the run used, version 2.1.4, which differs from the next release,
\texttt{v2.1.5}, only in the files that the latter's notes list, and the document as
shipped with that release: a template, whose hash line was blank. The copy that the run
checked differed from it only in that line; the configuration file was byte for byte
\path{configs/phase1_rec8_v2.json} of the version described here, whose canonical hash is
the one in that line of the copy that the run checked. A first attempt of the run, on 14
July 2026, stopped before any training, because the corpus that it had generated carried no
labels; the labels were computed separately on 15 July 2026, and the run was started again,
with the same code and configuration, on 16 July. The document of the joint-embedding run
was not in the repository before that run. On 2 August 2026 at 17:29 UTC the repository's
history, which until then consisted of the first commit alone, was replaced, by a force
push, with a new history whose first commit (\texttt{a548825}, tagged \texttt{v2.1.5}),
dated 17:26 UTC that day (3 August in the committer's time zone, UTC+8), contains both
documents with their hash lines filled and the repository's provenance note. The note,
compiled on 1 August 2026, records the hashes of the documents that the runs checked, which
are those of the committed documents, and the hashes of the reports and of the archives
that contained them. The first commit of the replaced history is kept under the tag
\texttt{provenance-2026-07-14}, created in October 2026. An export of the hosting service's
record of the repository's activity, which logs the creation of the branch at that commit
on 14 July and the force push of 2 August, and the security log's entries that name the
repository are committed with the records, as are the three reports used here, with their
hashes (Section~\ref{app:repro}).

\paragraph{Audit}
An independent script recomputed the three-dimensional verdict from the report and agreed
on \numAuditItems{} items. Its first pass reported one discrepancy, which was in the audit
script itself: it had combined the gate's conditions more strictly than the pre-registered
rule. Both outputs are archived.
